\section{Fundamental operators of the TPK: domains, actions, invariants, and memory}
\label{sec:inventario-operatorio-tpk-rev11}
\index{TPK!operational inventory}

The \cref{def:categorias-operatorias-tpk-rev11} first fixed the eight
ontological--operatorial classes: support and state; selection; transport;
readers and quotients; memory and boundary; periodicity and return;
prolongation; and outputs. This section defines their realizations. For each
operator, the sextuple \((D,C,f,I,L,M)\) of domain, codomain, rule,
invariants, published or lost information, and preserved memory is declared. The subsequent
nominal census groups symbols already defined by these six data and does not
replace their domains or rules.

\begin{definition}[Typed Record for a TPK Operator]
A canonical operator is a sextuple
\[
 \mathcal O=(D_{\mathcal O},C_{\mathcal O},f_{\mathcal O},
 I_{\mathcal O},L_{\mathcal O},M_{\mathcal O}),
\]
where \(f_{\mathcal O}:D_{\mathcal O}\to C_{\mathcal O}\) is the action rule, \(I_{\mathcal O}\) is the family of relations it preserves, \(L_{\mathcal O}\) declares the published coordinates and the equivalence relation induced by forgetting, and \(M_{\mathcal O}\) enumerates the data that remain in the enriched state after publication. Two realizations represent the same operator only when a typed equivalence exists that intertwines their rules and invariants; sharing a period, cardinality, or visible output is insufficient.
\end{definition}

\subsection{Modular readers built}

The nonadic reduction is fixed by the published representative.
\begin{equation}
 \rho_9(n)=1+((n-1)\bmod 9)
 =n-9\left\lfloor\frac{n-1}{9}\right\rfloor,
 \label{eq:rho9-construida-tpk-rev11}
\end{equation}
so that the null class is written as \(9\). For \(n\geq1\), the
digital root satisfies
\[
 \operatorname{dr}_9(n)=\rho_9(n),
\]
but the two expressions preserve distinct semantics: the first is a
section of the quotient; the second publishes the iterated sum of digits. The tritic
lift writes uniquely
\begin{equation}
 n=r_3(n)+3c_3(n),
 \qquad r_3(n)\in\{-1,0,+1\},\quad c_3(n)\in\mathbb Z,
 \label{eq:lift-mod3-tpk-rev11}
\end{equation}
and simultaneously preserves oriented residue and carry quotient.
Finally, the decimal-block reader is the Euclidean division
\begin{equation}
 N=1000q_{1000}(N)+r_{1000}(N),
 \qquad 0\leq r_{1000}(N)<1000,
 \label{eq:lector-mod1000-tpk-rev11}
\end{equation}
therefore basis \(1000\), residue modulo \(1000\) and carry \(q_{1000}\)
are three data related and not a same object.

\subsection{Classification of Operator Realizations by Domain and Function}

The preceding definitions govern the reading of the table. Its fourteen
functional groupings refine the eight operator classes and summarize the
realizations that share a function; a class may therefore occupy more
than one row. The table does not operate as a replacement definition. Symbols that
share a cardinal remain separated by domain and codomain.

\begingroup
\footnotesize
\renewcommand{\arraystretch}{1.06}
\begin{longtable}{@{}p{.21\textwidth}p{.31\textwidth}p{.40\textwidth}@{}}
\toprule
Functional grouping & Operators and Spaces & Domain, Action and Invariants \\
\midrule
\endfirsthead
\toprule
Grouping functional & Operators and Spaces & Domain, Action and Invariants \\
\midrule
\endhead
State and route memory &
\(\mathcal X_{\rm TPK}^{\rm enr}\), fiber \(\mathcal F_q\), route memory &
They preserve position, leaf, orientation, residue, quotient, carry, signature,
boundary, cylinder, prefix, genealogy, and memory. \\

Sheet selection &
\(M_{\rm ph}:\{1,\ldots,9\}\to\{-1,0,+1\}\) &
It activates the additive evaluation, the multiplicative evaluation, or the threshold passage; it does not
by itself move the state. \\

Transporte espacial &
\(T_N,T_E,T_S,T_O\), reversible two-cursor engine,
\(F_{\rm obs}\) &
It realizes the translations of the Cayley graph, updates sheet and orientation and
preserves each transition of the path. \\

Lectores modulares &
\(\rho_9\), \((r_3,c_3)\), \((q_{1000},r_{1000})\) &
Publish nonadic phase, tritic orientation and decimal block without erasing the lift quotients. \\

Block and memory &
\(\mathcal K_{27}=F_{\rm obs}^{27}\), memory filter
\(\mathcal N_{27}\) &
The first composes twenty-seven transports; the second selects the
memory events. Equal cardinality does not imply an equal operator. \\

Twelve-phase state &
\(\mathcal R_{12}\), centers \(\theta_m=30m-10\) degrees,
refinement \(\mathcal R_{120}\) &
Organize the twelve phases, their opposition, and the angular refinement in increments of
three degrees. \\

Periodicity and depths &
\(C_{12}\hookrightarrow C_{108}\twoheadrightarrow C_9\),
\(\Pi_{108}^{\rm cur}\), \(w_{108}\), \(w_{1080}\),
\(\Gamma_0(1080)\) &
It distinguishes the total periodicity of order \(108\), the cursor projection, the
words of 108 and 1080 trits, and modular level 1080. \\

Channels (90/120) &
mask \(C_m\), \(S_{90},S_{120}\), orbital grammar,
supertrace, and invariants \(\nu_{120},\nu_{270}\) &
It separates the two genealogies, their geometric tails, their weights, and the
coordinates that they extract from the state. \\

Bidirectionality &
\(P_\pm\), route inversion, \(\operatorname{Ext}_\pm\),
\(N_\pm\), \(\beta\), \(C_M=-\operatorname{rev}\) &
It implements, respectively, spectral exchange, spatial outward passage and return,
future/past prolongation, valence, balance, and word conjugation. \\

Emission and coinduction &
\(L_0,L_1\), \(R_{36}\), \(G_9\), partition \(3^6=729\),
\(\varprojlim\operatorname{Cyl}_m\) &
It lifts the seed, opens the boundary, and exhaustively enumerates the
finite prolongations. The numerical limit and the enriched lift
belong to distinct types. \\

Numeric Readers &
\(\mathscr C_\pi\), \(\mathscr P_e\), \(F_{\rm av}\) &
Through rational forks, factorial recurrence, and Fibonacci
self-scaling, they construct the published sections of \(\pi,e,\varphi\). \\

Exceptional Projection &
Witt vector lifting, weight enumerator, and non-selective shadow &
Transports the same state to ternary incidence without intervening in the Archimedean selection of digits. \\

Atlas and realization &
route signature, atlas (81/56/13/324), involution \(C_M\), electronic
center, and fiber operators &
Convert the route state into classifiers, characters, and operators before
applying any dimensional chart. \\

Finite validation &
exhaustive censuses, ablations, and causal independence &
It checks coverage and uniqueness at the computed depth, and verifies that
no target prefix intervenes in the generating map. \\
\bottomrule
\end{longtable}
\endgroup

\Needspace{46\baselineskip}
\subsection{Functional Index of the Canonical Operators}

The preceding classification organizes the realizations by function. The following table
gathers the operators after the categorical definition and the constructed
readers. Each row declares domain, codomain, rule, and invariants; the
symbols make it possible to recognize a single definition throughout its
different realizations without identifying operators of different types.

\begingroup
\footnotesize
\renewcommand{\arraystretch}{1.06}

\begin{longtable}{@{}r >{\raggedright\arraybackslash}p{.42\textwidth}
  >{\raggedright\arraybackslash}p{.22\textwidth}
  >{\raggedright\arraybackslash}p{.22\textwidth}@{}}
\toprule
\# & Canonical Operator & Symbol & Capa tipada \\
\midrule
\endfirsthead
\toprule
\# & Canonical Operator & Symbol & Capa tipada \\
\midrule
\endhead
1 & Additive--multiplicative arithmetic cell of order nine
  & \(\mathcal A_9\) & arithmetic support \\
2 & Residual decomposition and nonadic quotient
  & \((\rho_9,q_9)\) & reduction and elevation \\
3 & Ternary oriented reduction
  & \(\rho_3^{\rm or}\) & reduction and elevation \\
4 & Ternary-decimal positional reader of base thousand
  & \(\iota_{1000}\) & reduction and elevation \\
5 & Cardinal Action Oriented on the Arithmetic Torus
  & \(T_d,\ d\in\{N,E,S,O\}\) & transporte cardinal \\
6 & Additive--multiplicative local feedback
  & \(F_{\rm loc}\) & transporte cardinal \\
7 & Minimum emitter of two cursors
  & \(T_{\min}\) & finite emission \\
8 & Decimal Arithmetic Sign Coupler
  & \(G\) & finite emission \\
9 & Ternary projection of the catalog
  & \(\Psi_6\) & finite emission \\
10 & Local Surjectivity of Cardinal Transport
  & \(\mathcal R_3\) & transporte cardinal \\
11 & Transportation loop of twenty-seven iterations
  & \(\mathcal K_{27}\) & transporte cardinal \\
12 & Enriched state of the topological phase core
  & \(\mathcal X_{\rm TPK}^{\rm enr}\) & estado enriquecido \\
13 & Hensel State Transition
  & \(\mathcal H\) & estado enriquecido \\
14 & Exact Transition of Ternary-Decimal Cylinders
  & \(\mathcal C_{3,1000}\) & estado enriquecido \\
15 & Joint boundary signature
  & \(\Sigma_{\rm B}\) & estado enriquecido \\
16 & Relationship of extension and nonadic survival
  & \({\rm EF}\text{--}G_9\) & prolongation and monodromy \\
17 & Nonadic monodromy with boundary holonomy
  & \(\mathcal M_9\) & prolongation and monodromy \\
18 & Global Section of the Inverse Survivor System
  & \(X_\chi=\varprojlim_m{\rm Surv}^{(\chi)}_m\)
  & prolongation and monodromy \\
19 & Bivariant Category of Arithmetic Paths
  & \(\mathcal P_{\rm APP}^{(2)}\) & transporte cardinal \\
20 & Observable chronology of two cursors
  & \(F_{\rm obs}\) & transporte cardinal \\
21 & Decimal Alphabet of Forty-Two Trinities
  & \(\mathcal A_{42}\) & finite emission \\
22 & Typed extension relation of fibers
  & \({\rm Ext}_r\) & estado enriquecido \\
23 & Nonadic transfer of compatible fibers
  & \(K_{\rm ph}\) & prolongation and monodromy \\
24 & Refinement of Seeds by Advancement and Rotation
  & \((P,H,Q)\) & transporte cardinal \\
25 & Mode and Orientation Cocycles
  & \((c_{\rm mode},c_{\rm or})\) & estado enriquecido \\
26 & Triadic prolongation of twenty blocks
  & \(E_{20}\) & prolongation and monodromy \\
27 & Dodecaphonic oriented system
  & \(\mathcal R_{12}\) & dodecaphonic reading \\
28 & Signed harmonic dodecaphonic observable
  & \(U_{12}^{\rm sgn}\) & dodecaphonic reading \\
29 & Harmonic transformation by three-step orbits
  & \(\mathcal H_{12}\) & dodecaphonic reading \\
30 & Cyclic difference transversal extractor
  & \(\mathcal E_{3,4}\) & dodecaphonic reading \\
31 & Dodecaphasic closure recurrence in base thousand
  & \(\mathcal C_\alpha\) & close in base thousand \\
32 & Enhanced history transducer to angular character
  & \((L_{\rm tick},\chi_{\rm mass})\) & interfaz dimensional \\
33 & Bidiirectional lattice of index twelve
  & \(\mathcal W_{\pm}\) & interfaz dimensional \\*
34 & Nonadic Atlas of Extension Families
  & \(\mathcal A_{\rm species}\) & interfaz dimensional \\
\bottomrule
\end{longtable}
\endgroup

The symbol \(L_{\rm tick}\) denotes the iteration and return reader. Its
codomain is the temporal memory of the state; it does not define an independent physical
time unit.

The internal selector function
\[
 M_{\rm ph}:\mathcal D_9\longrightarrow\{+1,-1,0\}
\]
belongs to the TPK\@; its tabulated vector is
\[
 m_{\rm ph}:=\bigl(M_{\rm ph}(1),\ldots,M_{\rm ph}(9)\bigr)^{\mathsf T}.
\]
The function \(M_{\rm ph}\), the vector \(m_{\rm ph}\) and the state operator
are objects of different types. The scalar value \(M_{\rm ph}(g(t))\) enters
as the second argument of the selector \(\operatorname{Sel}_t\); it is not composed
directly with the state. Cardinal transport \(T_d\) occupies another factor
of the transition \(\mathcal U_t\) and performs movement on APP through
its lift \(\operatorname{Lift}(T_d)\). This separation preserves a single
container, distinguishes phase selection from transport, and preserves the
complete architecture by types.

\Needspace{9\baselineskip}
\begin{resultbox}[title={Finite depth and enriched limit}]
The enumeration, censuses, and contraction of numerical cylinders are
exact results in their domains. The existence and uniqueness of an enriched
section jointly subject to all survival predicates
follow when the lifting fibers are nonempty, singletons, and
compatible at every depth. A calculation on a finite window
verifies the fibers in that window, but does not by itself imply the
quantifier over all depths.
\end{resultbox}

\begin{theorem}[Composition of the internal dynamics]
\label{thm:composicion-dinamica-tpk-rev11}
A transition observable of the TPK is the composition typed
\[
 \operatorname{Tra}_t(y,d):=\operatorname{Lift}(T_d)(y),
 \qquad
 \operatorname{Upd}_t:=\operatorname{Rec}_t\circ\operatorname{Mem}_t,
\]
\[
 \mathcal U_t(x)
 =\operatorname{Upd}_t\!\left(
  \operatorname{Tra}_t\!\left(
   \operatorname{Sel}_t\!\left(x,M_{\rm ph}(g(t))\right),d(t)
  \right)\right)
\]
on \(\mathcal X_{\rm TPK}^{\rm enr}\), followed, when appropriate, by
a modular reader, a cylindrical extension, or one of the two projections.
Each factor acts on a distinct component of the state; for that reason the
selector, the movement, the clock, the memory, and the publication cannot
collapse into a single modular reduction.
\end{theorem}

The preceding table enumerates the operator families. Equivalences among
local realizations must intertwine their domains, rules, and invariants.
Subsequent developments select compositions of these operations;
they introduce no fourth primitive object alongside APP, TRIT, and TPK.
